\documentclass[11pt,twoside]{article}

\usepackage[T1]{fontenc}
\usepackage{newpxtext}
\usepackage{amsmath,amsthm,mathtools}
\usepackage{newpxmath}
\usepackage{booktabs}
\usepackage[
  activate={true,nocompatibility},
  tracking=true,
  expansion=true,
  factor=1100,
  stretch=10,
  shrink=10
]{microtype}
\microtypecontext{spacing=nonfrench}
\usepackage{geometry}
\usepackage[dvipsnames]{xcolor}
\usepackage{fancyhdr}
\usepackage{titlesec}
\usepackage{setspace}
\usepackage{needspace}
\usepackage{etoolbox}
\usepackage[most]{tcolorbox}

\geometry{
  paperwidth=8.5in,
  paperheight=11in,
  left=1.2in,
  right=1.2in,
  top=1.2in,
  bottom=1.2in,
  headheight=15pt
}

\widowpenalty=10000
\clubpenalty=10000
\raggedbottom

\setlength{\abovedisplayskip}{12pt plus 3pt minus 4pt}
\setlength{\belowdisplayskip}{10pt plus 3pt minus 4pt}
\setlength{\abovedisplayshortskip}{6pt plus 3pt}
\setlength{\belowdisplayshortskip}{8pt plus 3pt minus 3pt}

\definecolor{sectioncolor}{RGB}{25,25,45}
\definecolor{linkcolor}{RGB}{60,50,130}
\definecolor{citecolor}{RGB}{100,50,50}
\definecolor{theorembar}{RGB}{50,50,80}
\definecolor{defnbg}{RGB}{245,243,248}
\definecolor{remarkbar}{RGB}{140,140,160}

\titleformat{\section}
  {\normalfont\Large\bfseries\color{sectioncolor}}
  {\thesection}{1em}{}
\titleformat{\subsection}
  {\normalfont\large\bfseries\color{sectioncolor}}
  {\thesubsection}{1em}{}
\titleformat{\subsubsection}
  {\normalfont\normalsize\bfseries\color{sectioncolor}}
  {\thesubsubsection}{1em}{}

\titlespacing*{\section}{0pt}{3.5ex plus 1.5ex minus 1ex}{2ex plus .2ex}
\titlespacing*{\subsection}{0pt}{2.5ex plus 1ex minus .5ex}{1.2ex plus .2ex}
\titlespacing*{\subsubsection}{0pt}{2ex plus 0.8ex minus .3ex}{1ex plus .2ex}

\pretocmd{\section}{\needspace{6\baselineskip}}{}{}
\pretocmd{\subsection}{\needspace{5\baselineskip}}{}{}

\pagestyle{fancy}
\fancyhf{}
\fancyhead[LE,RO]{\small\thepage}
\fancyhead[RE]{\small\itshape Microprice and Semi-Markov Reinforcement Learning}
\fancyhead[LO]{\small\itshape L.\,Petersen}
\renewcommand{\headrulewidth}{0.3pt}
\renewcommand{\footrulewidth}{0pt}

\fancypagestyle{plain}{
  \fancyhf{}
  \fancyfoot[C]{\small\thepage}
  \renewcommand{\headrulewidth}{0pt}
}

\theoremstyle{plain}
\newtheorem{theorem}{Theorem}[section]
\newtheorem{proposition}[theorem]{Proposition}
\newtheorem{lemma}[theorem]{Lemma}
\newtheorem{corollary}[theorem]{Corollary}

\theoremstyle{definition}
\newtheorem{definition}[theorem]{Definition}

\theoremstyle{remark}
\newtheorem{remark}[theorem]{Remark}

\tcolorboxenvironment{theorem}{
  blanker, borderline west={2.5pt}{0pt}{theorembar},
  left=10pt, before skip=10pt plus 2pt minus 2pt,
  after skip=10pt plus 2pt minus 2pt, breakable
}
\tcolorboxenvironment{proposition}{
  blanker, borderline west={2.5pt}{0pt}{theorembar},
  left=10pt, before skip=10pt plus 2pt minus 2pt,
  after skip=10pt plus 2pt minus 2pt, breakable
}
\tcolorboxenvironment{lemma}{
  blanker, borderline west={2.5pt}{0pt}{theorembar},
  left=10pt, before skip=10pt plus 2pt minus 2pt,
  after skip=10pt plus 2pt minus 2pt, breakable
}
\tcolorboxenvironment{corollary}{
  blanker, borderline west={2.5pt}{0pt}{theorembar},
  left=10pt, before skip=10pt plus 2pt minus 2pt,
  after skip=10pt plus 2pt minus 2pt, breakable
}
\tcolorboxenvironment{definition}{
  colback=defnbg, colframe=defnbg, boxrule=0pt, arc=2pt,
  left=8pt, right=8pt, top=6pt, bottom=6pt,
  before skip=10pt plus 2pt minus 2pt,
  after skip=10pt plus 2pt minus 2pt, breakable
}
\tcolorboxenvironment{remark}{
  blanker, borderline west={1.5pt}{0pt}{remarkbar},
  left=10pt, before skip=8pt plus 2pt minus 2pt,
  after skip=8pt plus 2pt minus 2pt, breakable
}

\AtEndEnvironment{proof}{\medskip}

\newcommand{\1}{\mathbf{1}}
\newcommand{\R}{\mathbb{R}}
\newcommand{\E}{\mathbb{E}}
\newcommand{\Prob}{\mathbb{P}}
\newcommand{\norm}[1]{\lVert #1 \rVert}
\newcommand{\PiB}{\Pi_B}
\newcommand{\PiP}{\Pi_P}

\usepackage{hyperref}
\hypersetup{
  colorlinks=true,
  linkcolor=linkcolor,
  citecolor=citecolor,
  urlcolor=linkcolor,
  bookmarksnumbered=true,
  pdfauthor={Lev Petersen},
  pdftitle={Microprice, Bellman Values, and Semi-Markov Reinforcement Learning}
}

\input{rl_note_summary_macros.tex}

\begin{document}

\begin{titlepage}
\centering

\vspace*{1.4cm}

{\Huge\bfseries\color{sectioncolor} Microprice, Bellman Values,\\[0.5em]
Semi-Markov Reinforcement Learning\par}

\vspace{0.45cm}

{\large\color{sectioncolor} Paper II: Bellman values and learning extensions of\\[0.3em]
\textit{Microprice as a Centered Poisson Corrector}\par}

\vspace{0.4cm}
{\color{sectioncolor}\rule{0.52\textwidth}{0.8pt}}

\vspace{1.0cm}

{\Large Lev Petersen\par}

\vspace{1.3cm}

\begin{minipage}{0.86\textwidth}
\setstretch{1.14}
\small
\noindent\textbf{Abstract.}
This companion note rewrites the move-time objects of \cite{Petersen2025} in value-function language and adds exact event-time targets together with online learning rules on the same price-move episodes. Let
\[
r(x):=\E[p_{n+1}-p_n\mid X_n=x],
\qquad
N:=(I-Q)^{-1},
\qquad
G_1:=Nr,
\qquad
B:=NT.
\]
If $B$ is ergodic with stationary distribution $\pi$, set
\[
\mu=\pi^\top G_1,
\qquad
\PiB:=\1\pi^\top,
\qquad
g=(I-B+\PiB)^{-1}(G_1-\mu\1),
\qquad
G^*:=(I-B+\PiB)^{-1}G_1=g+\mu\1.
\]
Then classical microprice is already the move-time relative value function,
\[
g=G_1-\mu\,\1+Bg,
\qquad
\pi^\top g=0.
\]
The event-time discounted Bellman value
\[
V_\gamma=(I-\gamma P)^{-1}r
\]
factorizes as
\[
V_\gamma=G_{1,\gamma}+B_\gamma V_\gamma,
\qquad
G_{1,\gamma}=(I-\gamma Q)^{-1}r,
\qquad
B_\gamma=\gamma(I-\gamma Q)^{-1}T.
\]
We reserve $\gamma$ for event-time discounting; the move-time resolvent $(I-\beta B)^{-1}G_1$ of Paper~I discounts only after the reduction. For deterministic horizons,
\[
V^{(H)}=\sum_{k=0}^{H-1}P^k r,
\qquad
V^{(H)}=H\mu_P\,\1+g_P-P^H g_P,
\]
where
\[
g_P=(I-P+\PiP)^{-1}(r-\mu_P\1),
\qquad
\mu_P=\pi_P^\top r.
\]
Thus $g_P$ is the event-time analogue of the move-time boundary object $G^*$.

On the bundled public panels (\RLEmpiricalCellCount\ dataset-horizon cells), the fixed-horizon Bellman value beats classical microprice in \RLEmpiricalFiniteWins/\RLEmpiricalCellCount\ cells with geometric-mean grouped-RMSE ratio \RLEmpiricalFiniteGeo, and the horizon-matched discounted value wins in \RLEmpiricalDiscountWins/\RLEmpiricalCellCount\ cells with ratio \RLEmpiricalDiscountGeo. The event-time boundary value $g_P$ coincides with $G^*$ on the centered Stoikov and WSELOB panels, improves all six IEX cells with geometric-mean ratio \RLEmpiricalDiffIEXGeo, and the move-time formulation compresses event-time TD updates by a geometric-mean factor \RLEmpiricalMoveCompressionGeo{}x (median \RLEmpiricalMoveCompressionMedian{}x).

\vspace{0.9em}
\noindent\textbf{Keywords:} microprice, reinforcement learning, semi-Markov decision process, temporal-difference learning, value function, successor representation, average reward
\end{minipage}

\vfill

{\small\color{sectioncolor} Draft \textbullet\ \today}

\end{titlepage}

% ============================
\section{Introduction}

The main paper \cite{Petersen2025} identifies classical microprice as a centered Poisson corrector on the move-time chain. This note keeps the same reduction and asks two Bellman questions: what are the exact event-time forecasting targets on the full chain $P$, and how do those targets induce online learning rules on the same price-move episodes?

We therefore keep the same core dictionary:
\[
r(x):=\E[p_{n+1}-p_n\mid X_n=x],
\qquad
N:=(I-Q)^{-1},
\qquad
G_1:=Nr,
\qquad
B:=NT.
\]
If the move-time chain $B$ is ergodic with stationary distribution $\pi$, we write
\[
\mu:=\pi^\top G_1,
\qquad
\PiB:=\1\pi^\top,
\qquad
(I-B)g=G_1-\mu\1,
\qquad
\pi^\top g=0,
\]
and
\[
G^*:=(I-B+\PiB)^{-1}G_1=g+\mu\1.
\]
Paper~I focuses on the move-time Poisson representation, the regularizer $(I-\beta B)^{-1}G_1$, and the horizon effect. Paper~II adds the event-time Bellman values $V^{(H)}$ and $V_\gamma$, the event-time differential value $g_P$, and the corresponding TD/SMDP learning rules.

\begin{table}[t]
\centering
\small
\begin{tabular}{@{}p{3.6cm}p{8.1cm}@{}}
\toprule
Object & Role in the two-note sequence \\
\midrule
$r$, $N$, $G_1$, $B$ & Common primitives from the move-time reduction \\
$g$, $G^*$ & Centered and uncentered move-time boundary objects from Paper~I \\
$V^{(H)}$, $V_\gamma$ & Exact event-time forecasting targets added here \\
$g_P$ & Event-time analogue of the move-time centered correction \\
$(I-\beta B)^{-1}G_1$ vs.\ $V_\gamma$ & Move-time discount after reduction vs.\ event-time discount before reduction \\
\bottomrule
\end{tabular}
\caption{Shared notation and division of labor between the two companion notes.}
\label{tab:paper_map}
\end{table}

Empirically, $V^{(H)}$ is the strongest forecasting object on the bundled public panels, $V_\gamma$ is the cleanest one-shot surrogate, and the move-time updates are the most ergonomic online learners.

% ============================
\section{Discounted move-time Bellman equations}

Let $X_n$ be the event-time Markov chain of top-of-book states and let $p_n$ be the mid-price. As in \cite{Petersen2025}, write
\[
P=Q+T,
\qquad
r(x):=\E[p_{n+1}-p_n\mid X_n=x].
\]
Here $Q$ collects no-move transitions and $T$ move transitions. Assume $\rho(Q)<1$ so that the fundamental matrix
\[
N:=(I-Q)^{-1}
\]
exists.

\begin{definition}[Move-time objects from Paper~I]
The one-move reward and move-time kernel are
\[
G_1:=Nr,
\qquad
B:=NT.
\]
These are exactly the objects appearing in \cite{Petersen2025}.
\end{definition}

For $\gamma\in(0,1)$ define the discounted event-time Bellman value
\begin{equation}\label{eq:discounted_event_value}
V_\gamma(x):=\E_x\!\left[\sum_{n=0}^{\infty}\gamma^n(p_{n+1}-p_n)\right].
\end{equation}
The first proposition shows that $V_\gamma$ can be written using the same move-time reduction as microprice.

\begin{proposition}[Discounted move-time factorization]\label{prop:discounted_move_factorization}
For $\gamma\in(0,1)$ define
\[
G_{1,\gamma}:=(I-\gamma Q)^{-1}r,
\qquad
B_\gamma:=\gamma(I-\gamma Q)^{-1}T.
\]
Then $V_\gamma$ satisfies the semi-Markov Bellman equation
\begin{equation}\label{eq:discounted_move_bellman}
V_\gamma = G_{1,\gamma} + B_\gamma V_\gamma.
\end{equation}
Equivalently,
\begin{equation}\label{eq:discounted_move_resolvent}
V_\gamma=(I-B_\gamma)^{-1}G_{1,\gamma}.
\end{equation}
Moreover,
\begin{equation}\label{eq:discounted_move_identity}
I-B_\gamma=(I-\gamma Q)^{-1}(I-\gamma P),
\end{equation}
so \eqref{eq:discounted_move_resolvent} is exactly the same value as
\begin{equation}\label{eq:discounted_event_resolvent}
V_\gamma=(I-\gamma P)^{-1}r.
\end{equation}
\end{proposition}

Equation \eqref{eq:discounted_move_bellman} is the most direct RL bridge from the original paper: the only change from $(G_1,B)$ is that each no-move excursion is discounted before it is collapsed.

\begin{remark}[Event-time discounting versus move-time discounting]
The discounted resolvent $(I-\beta B)^{-1}G_1$ studied in \cite{Petersen2025} discounts after the reduction to move time, so each price-move episode receives the same geometric weight $\beta^k$ regardless of its duration in event time. By contrast, $V_\gamma=(I-\gamma P)^{-1}r$ discounts before the reduction, and the induced pair $(G_{1,\gamma},B_\gamma)$ therefore carries episode-length-dependent discount factors. The two formulas agree only when every move occurs after one event.
\end{remark}

\begin{proposition}[Classical microprice as a move-time relative value]\label{prop:move_time_relative_value}
Assume the move-time chain $B$ is ergodic with stationary distribution $\pi$ and set
\[
\mu:=\pi^\top G_1,
\qquad
\PiB:=\1\pi^\top.
\]
Then the centered correction
\begin{equation}\label{eq:move_time_relative_value}
g:=(I-B+\PiB)^{-1}(G_1-\mu\1)
\end{equation}
satisfies
\begin{equation}\label{eq:move_time_average_reward_bellman}
g=G_1-\mu\1+Bg,
\qquad
\pi^\top g=0.
\end{equation}
Equivalently,
\[
G^*:=(I-B+\PiB)^{-1}G_1=g+\mu\1.
\]
So the centered microprice correction of \cite{Petersen2025} is precisely the average-reward relative value function of the move-time chain.
\end{proposition}

Thus the undiscounted microprice correction is already a value function; the discounted semi-Markov family simply replaces the average-reward boundary problem by a horizon-aware Bellman objective.

% ============================
\section{Event-time Bellman values and differential limits}

The event-time chain admits an exact finite-horizon Bellman family.

\begin{proposition}[Finite-horizon Bellman values]\label{prop:finite_horizon}
For each integer $H\ge 0$, define
\[
V^{(0)}:=0,
\qquad
V^{(h+1)}:=r+PV^{(h)}.
\]
Then
\begin{equation}\label{eq:finite_horizon_closed}
V^{(H)}=\sum_{k=0}^{H-1}P^k r,
\end{equation}
and, if $R_n:=p_{n+1}-p_n$,
\begin{equation}\label{eq:finite_horizon_expectation}
V^{(H)}(x)=\E_x\!\left[\sum_{n=0}^{H-1}R_n\right].
\end{equation}
Thus $V^{(H)}$ is the exact deterministic-horizon predictor of future mid-price change.
\end{proposition}

When $P$ is ergodic, the same Poisson-corrector structure that appears in \cite{Petersen2025} reappears on event time.

\begin{proposition}[Event-time differential value]\label{prop:event_time_poisson}
Assume $P$ is ergodic with stationary distribution $\pi_P$ and set
\[
\PiP:=\1\pi_P^\top,
\qquad
\mu_P:=\pi_P^\top r.
\]
Define the event-time differential value
\begin{equation}\label{eq:event_time_diff_value}
g_P:=(I-P+\PiP)^{-1}(r-\mu_P\1),
\qquad
\pi_P^\top g_P=0.
\end{equation}
Then, for every integer $H\ge 0$,
\begin{equation}\label{eq:finite_horizon_decomp}
V^{(H)}=H\mu_P\,\1+g_P-P^H g_P.
\end{equation}
For every $\gamma\in(0,1)$,
\begin{equation}\label{eq:discounted_split_event_time}
V_\gamma
=
\frac{\mu_P}{1-\gamma}\,\1
+
(I-\gamma P)^{-1}(I-\PiP)r.
\end{equation}
Hence
\begin{equation}\label{eq:discounted_boundary_event_time}
V_\gamma-\frac{\mu_P}{1-\gamma}\,\1
\longrightarrow g_P
\qquad\text{as }\gamma\uparrow 1.
\end{equation}
\end{proposition}

This is the exact event-time analogue of the main theorem in \cite{Petersen2025}: $G^*$ is the centered long-run value after the $Q/T$ reduction, while $g_P$ is the centered long-run value before that reduction.
% ============================
\section{Learning rules}

\subsection{Event-time TD and LSTD}

The discounted Bellman equation on event time,
\[
V_\gamma=r+\gamma PV_\gamma,
\]
may be learned online by standard temporal-difference learning \cite{Sutton1988}. With $X_t$ the discretized state and $R_t:=p_{t+1}-p_t$,
\begin{equation}\label{eq:td0}
V_{t+1}(X_t)
\leftarrow
V_t(X_t)+\eta_t\bigl(R_t+\gamma V_t(X_{t+1})-V_t(X_t)\bigr).
\end{equation}
With linear features $\phi(x)$, least-squares TD solves
\begin{equation}\label{eq:lstd}
A_T w_T=b_T,
\qquad
A_T:=\sum_{t<T}\phi_t(\phi_t-\gamma\phi_{t+1})^\top,
\qquad
b_T:=\sum_{t<T}\phi_t R_t,
\end{equation}
the standard linear policy-evaluation analogue \cite{BradtkeBarto1996}. The same value also admits the successor-representation factorization \cite{Dayan1993}
\begin{equation}\label{eq:sr}
M_\gamma:=(I-\gamma P)^{-1}=
\sum_{k=0}^{\infty}\gamma^k P^k,
\qquad
V_\gamma=M_\gamma r.
\end{equation}
If rewards are parameterized as $r_\theta=\Phi\theta$, then the successor features are $\Psi_\gamma:=M_\gamma\Phi$ and $V_{\gamma,\theta}=\Psi_\gamma\theta$ \cite{Barreto2018}.

\subsection{Move-time differential TD for classical microprice}

Proposition~\ref{prop:move_time_relative_value} shows that the centered correction of \cite{Petersen2025} is already an average-reward Bellman value on move time. Writing $\Delta P_k:=p_{\tau_{k+1}}-p_{\tau_k}$ for the move increment,
\begin{equation}\label{eq:move_time_relative_expectation}
g(y)=\E\bigl[\Delta P_k-\mu+g(Y_{k+1})\mid Y_k=y\bigr].
\end{equation}
Hence the original microprice correction may be learned online by differential TD. A tabular recursion is
\begin{equation}\label{eq:move_time_differential_td}
\delta_k:=\Delta P_k-\bar\mu_k+h_k(Y_{k+1})-h_k(Y_k),
\end{equation}
\begin{equation}\label{eq:move_time_differential_td_updates}
h_{k+1}(Y_k)\leftarrow h_k(Y_k)+\alpha_k\delta_k,
\qquad
\bar\mu_{k+1}\leftarrow \bar\mu_k+\beta_k\delta_k.
\end{equation}
This is the exact RL version of classical microprice: same move-time chain, same state space, same centered boundary object.

\subsection{Move-time semi-Markov TD for discounted forecasting}

For horizon-aware forecasting, replace the average-reward boundary problem by the discounted semi-Markov value. Let $\tau_k$ be the successive mid-price move times, $Y_k:=X_{\tau_k}$ the move-time chain, and $D_k:=\tau_{k+1}-\tau_k$ the move duration. Then the discounted value at move times satisfies the recursion
\begin{equation}\label{eq:smdp_bellman}
W_\gamma(y)=\E\bigl[\gamma^{D_k-1}\Delta P_k+\gamma^{D_k}W_\gamma(Y_{k+1})\mid Y_k=y\bigr].
\end{equation}
The factor $\gamma^{D_k-1}$ on $\Delta P_k$ appears because the price change is realized on the last event of the episode. The corresponding tabular SMDP-TD(0) update is
\begin{equation}\label{eq:smdp_td}
W_{k+1}(Y_k)
\leftarrow
W_k(Y_k)
+
\eta_k\bigl(\gamma^{D_k-1}\Delta P_k+\gamma^{D_k}W_k(Y_{k+1})-W_k(Y_k)\bigr).
\end{equation}
This uses the same move-time episodes as differential TD, but targets the discounted Bellman value rather than the centered long-run corrector. The choice is therefore objective-driven: differential TD learns the classical microprice boundary object, while SMDP-TD learns the horizon-aware discounted predictor. The model-based counterparts are the exact finite-horizon value $V^{(H)}$ and the one-shot surrogate $V_{\gamma_H}$ with $\gamma_H=1-1/H$.

% ============================
\section{Empirical results}

We evaluate the RL extension on the bundled public panels: Stoikov BAC/CVX, IEX BAC/SPY/GLD, and WSELOB PEKAO/PZU/PKOBP/PKNORLEN/KGHM. The grouped conditional-curve metric is the same one used by the existing public scripts.

\begin{table}[t]
\centering
\small
\begin{tabular}{@{}p{4.3cm}ccp{3.2cm}@{}}
\toprule
Predictor & Wins vs. $G^*$ & Geo.-mean RMSE ratio & Interpretation \\
\midrule
Finite-horizon Bellman $V^{(H)}$ & \RLEmpiricalFiniteWins/\RLEmpiricalCellCount & \RLEmpiricalFiniteGeo & exact fixed-horizon target \\
Discounted Bellman $V_{\gamma_H}$ & \RLEmpiricalDiscountWins/\RLEmpiricalCellCount & \RLEmpiricalDiscountGeo & one discounted solve \\
Event-time differential value $g_P$ & IEX 6/6; elsewhere ties & \RLEmpiricalDiffGeo & event-time boundary object \\
\bottomrule
\end{tabular}
\caption{Bundled-data comparison against classical microprice. Ratios below one favor the RL extension. For $g_P$, the non-IEX panels are equal to $G^*$ up to numerical tolerance, so the only visible differences come from the IEX stress cases.}
\label{tab:rl_summary}
\end{table}

Table~\ref{tab:rl_summary} gives the main forecasting message. The exact fixed-horizon Bellman value beats classical microprice in \RLEmpiricalFiniteWins/\RLEmpiricalCellCount\ cells, with geometric-mean RMSE ratio \RLEmpiricalFiniteGeo. The horizon-matched discounted value, with $\gamma_H=1-1/H$, wins in \RLEmpiricalDiscountWins/\RLEmpiricalCellCount\ cells with ratio \RLEmpiricalDiscountGeo. The largest gains are on short-horizon WSELOB cells; SPY at sixty steps remains the main negative case.

The event-time boundary value $g_P$ does exactly what the theory suggests. On the Stoikov and WSELOB panels, which are already close to the centered regime of \cite{Petersen2025}, $g_P$ and $G^*$ coincide to numerical tolerance. On the six IEX cells, where event time and move time are more visibly separated, $g_P$ improves all six cells with geometric-mean ratio \RLEmpiricalDiffIEXGeo. It is therefore the correct event-time boundary object, though not a replacement for the objective-matched finite-horizon value.

\begin{table}[t]
\centering
\small
\begin{tabular}{lcc}
\toprule
Panel & Geo.-mean events per move & Interpretation \\
\midrule
Stoikov & \RLEmpiricalMoveCompressionStoikovGeo{}x & modest temporal aggregation \\
IEX & \RLEmpiricalMoveCompressionIEXGeo{}x & moderate update compression \\
WSELOB & \RLEmpiricalMoveCompressionWSEGeo{}x & very sparse move times \\
Overall & \RLEmpiricalMoveCompressionGeo{}x & median \RLEmpiricalMoveCompressionMedian{}x \\
\bottomrule
\end{tabular}
\caption{Online-learning compression from move-time TD updates. One move-time update --- differential TD for classical microprice or SMDP-TD for discounted forecasting --- replaces one event-time TD step for each no-move event absorbed into the current price-move episode.}
\label{tab:compression}
\end{table}

Table~\ref{tab:compression} records the ergonomic gain behind the move-time RL formulation. Both the differential-TD microprice learner and the discounted SMDP learner update only once per price move, so the update count is compressed by a geometric-mean factor \RLEmpiricalMoveCompressionGeo{}x overall, with median \RLEmpiricalMoveCompressionMedian{}x and range \RLEmpiricalMoveCompressionMin{}x to \RLEmpiricalMoveCompressionMax{}x. This matters most on the sparse-move panels, where WSELOB averages \RLEmpiricalMoveCompressionWSEGeo{}x and PEKAO reaches \RLEmpiricalMoveCompressionMax{}x.

The discounted linear solve remains the cleanest model-based object. On the bundled models, solving $(I-\gamma P)V=r$ is \RLEmpiricalDiscountSpeedMean{}x cheaper on average than the classical microprice solve stage. In practice, $V^{(H)}$ is the exact forecasting target, $V_{\gamma_H}$ the simplest fast surrogate, and move-time SMDP the most ergonomic online learner.

% ============================
\section{Conclusion}

The reinforcement-learning extension of microprice is best read as Paper~II of the same operator story. Paper~I reduces event time to move time and identifies $G^*$ as the long-run move-time boundary object; the present note keeps that reduction fixed and adds the exact event-time Bellman objectives $V^{(H)}$, $V_\gamma$, and $g_P$.

This keeps the mathematics and the implementation aligned. The exact fixed-horizon Bellman value $V^{(H)}$ matches the forecasting objective, $V_\gamma$ gives the cleanest one-shot solve, $g_P$ is the correct event-time boundary object, the move-time differential recursion learns the classical microprice correction itself, and the move-time SMDP recursion learns the horizon-aware discounted predictor on the same price-move episodes. Empirically, fixed-horizon Bellman values deliver the largest forecasting gains, while move-time RL updates deliver the largest ergonomic gain.

\appendix

% ============================
\section{Proofs}

\begin{proof}[Proof of Proposition~\ref{prop:discounted_move_factorization}]
Partition the discounted return at the first price move. Summing over $n\ge 0$ no-move steps before that move gives
\[
V_\gamma
=
\sum_{n\ge 0}\gamma^n Q^n r
+
\sum_{n\ge 0}\gamma^{n+1}Q^n T V_\gamma
=
(I-\gamma Q)^{-1}r+\gamma(I-\gamma Q)^{-1}T V_\gamma,
\]
which is \eqref{eq:discounted_move_bellman}. Rearranging yields \eqref{eq:discounted_move_resolvent}. Also,
\[
(I-\gamma Q)(I-B_\gamma)=I-\gamma Q-\gamma T=I-\gamma P,
\]
which is \eqref{eq:discounted_move_identity}; multiplying \eqref{eq:discounted_move_resolvent} by $I-\gamma P$ gives \eqref{eq:discounted_event_resolvent}.
\end{proof}

\begin{proof}[Proof of Proposition~\ref{prop:move_time_relative_value}]
By definition,
\[
(I-B+\PiB)g=G_1-\mu\1.
\]
Since $\PiB g=\1\pi^\top g=0$, this reduces to
\[
(I-B)g=G_1-\mu\1,
\]
which is \eqref{eq:move_time_average_reward_bellman}. Conversely, \eqref{eq:move_time_average_reward_bellman} together with $\pi^\top g=0$ implies $\PiB g=0$ and therefore \eqref{eq:move_time_relative_value}.
\end{proof}

\begin{proof}[Proof of Proposition~\ref{prop:finite_horizon}]
The recursion telescopes to
\[
V^{(H)}=r+Pr+\cdots+P^{H-1}r=\sum_{k=0}^{H-1}P^k r,
\]
which is \eqref{eq:finite_horizon_closed}. Since $P^k r(x)=\E_x[R_k]$, summing over $k=0,\dots,H-1$ gives \eqref{eq:finite_horizon_expectation}.
\end{proof}

\begin{proof}[Proof of Proposition~\ref{prop:event_time_poisson}]
Since $g_P$ solves $(I-P)g_P=r-\mu_P\1$, we have $r=\mu_P\1+g_P-Pg_P$. Hence
\[
V^{(H)}
=
\sum_{k=0}^{H-1}P^k r
=
\sum_{k=0}^{H-1}P^k(\mu_P\1+g_P-Pg_P)
=
H\mu_P\1+g_P-P^H g_P,
\]
which is \eqref{eq:finite_horizon_decomp}. For the discounted identity, decompose
\[
r=\PiP r+(I-\PiP)r=\mu_P\1+(I-\PiP)r.
\]
Because $P\1=\1$ and $P\PiP=\PiP$,
\[
(I-\gamma P)^{-1}\PiP r=\sum_{k=0}^{\infty}\gamma^k\PiP r=\frac{\mu_P}{1-\gamma}\,\1.
\]
This gives \eqref{eq:discounted_split_event_time}. On the centered subspace $(I-\PiP)\R^L$, $P^k\to 0$, so the Neumann series converges to $(I-P+\PiP)^{-1}(I-\PiP)r=g_P$, yielding \eqref{eq:discounted_boundary_event_time}.
\end{proof}

\newpage
\begingroup
\small
\setlength{\itemsep}{3pt plus 1pt minus 1pt}
\setlength{\parsep}{0pt}
\begin{thebibliography}{99}

\bibitem{Petersen2025}
L.~Petersen, ``Microprice as a Centered Poisson Corrector,'' preprint, 2025.

\bibitem{Sutton1988}
R.~S.~Sutton, ``Learning to Predict by the Methods of Temporal Differences,'' \textit{Machine Learning}, vol.~3, no.~1, pp.~9--44, 1988. \href{https://doi.org/10.1023/A:1022633531479}{doi:10.1023/A:1022633531479}.

\bibitem{BradtkeBarto1996}
S.~J.~Bradtke and A.~G.~Barto, ``Linear Least-Squares Algorithms for Temporal Difference Learning,'' \textit{Machine Learning}, vol.~22, pp.~33--57, 1996. \href{https://doi.org/10.1023/A:1018056104778}{doi:10.1023/A:1018056104778}.

\bibitem{Dayan1993}
P.~Dayan, ``Improving Generalization for Temporal Difference Learning: The Successor Representation,'' \textit{Neural Computation}, vol.~5, no.~4, pp.~613--624, 1993. \href{https://doi.org/10.1162/neco.1993.5.4.613}{doi:10.1162/neco.1993.5.4.613}.

\bibitem{Barreto2018}
A.~Barreto, D.~Borsa, J.~Quan, T.~Schaul, D.~Silver, M.~Hessel, D.~Mankowitz, A.~Zidek, and R.~Munos, ``Transfer in Deep Reinforcement Learning Using Successor Features and Generalised Policy Improvement,'' in \textit{Proceedings of the 35th International Conference on Machine Learning}, PMLR~80, pp.~501--510, 2018.

\end{thebibliography}
\endgroup

\end{document}
